% !TEX root = ../../技术文档.tex
\section{心得体会}

本次参赛，我们以自主设计 RISC-V 处理器为核心，希望不仅完成 CPU 的功能验证和性能测试，还能将其部署到国产 FPGA 上，通过实际视觉应用检验处理器与硬件加速器的协同工作能力。在研发过程中，我们逐渐认识到，单个模块的性能固然重要，但如何在有限的硬件资源下完成系统集成，并保证各模块稳定运行，同样是 FPGA 设计中不可忽视的挑战。

CPU设计是本次作品投入精力最多的部分。我们从已有的五级流水线出发，逐步完成六级流水线、动态分支预测、Cache 和可配置双发射结构的设计。双发射的实现尤其让我们体会到处理器设计的复杂性：增加第二条执行通路只是其中一步，还必须解决指令配对、数据相关、结果前递、共享运算资源和异常恢复等问题。为了保证执行正确性，我们针对分支恢复、访存等待和错误路径副作用设计了定向测试。在一次零比较分支配对测试中，我们还发现测试平台自身的输入条件与现有设计不一致。这让我们认识到，验证不仅是检查设计能否通过测试，还需要保证测试条件与硬件实际行为相符。

NPU 与视频系统的集成使我们对软硬件协同设计有了更加具体的理解。 NPU 的设计需要从 YuNet 的实际算子和数据规模出发，在计算并行度、片上缓冲容量和数据搬运开销之间作出取舍。视频链路则不仅需要完成采集和图像叠加，还必须确保检测结果与原始画面对应。为此，我们设计了三帧槽管理和显示释放确认机制，并通过流式合成实现检测框及透明头像叠加。在头像功能中，利用无损行区间缓存减少片上纹理存储占用，也使我们更加直观地认识到，针对具体应用特点优化存储组织，往往比单纯增加硬件资源更有效。

在 FPGA 上板过程中，我们进一步体会到，仿真正确并不意味着系统集成就已经完成。时钟域交互、DDR 访问、存储器响应以及布局布线后的时序约束，都可能成为影响系统运行的关键因素。因此，我们采用模块验证、系统联合仿真与板级测试相结合的方式，逐步排查问题，并通过统一接口和配置管理减少不同模块集成时的差异。

目前，作品已经实现完整的人脸检测与头像叠加板级演示，但整网推理速度仍有较大的优化空间。这也让我们认识到，硬件设计不能只关注理论计算峰值，还应当结合实际执行周期，分析数据搬运、控制调度及计算单元的利用情况。

通过本次比赛，我们不仅加深了对 RISC-V 处理器微架构和 FPGA 硬件设计的理解，也积累了从 RTL 开发、功能验证到系统集成和实板调试的工程经验。相比单独完成某个模块，能够让自主设计的 CPU、NPU 和视频系统共同运行，并通过真实应用检验设计，是本次参赛最有价值的收获。
